\begin{itemize}
\item \textbf{Stylized decision layer.} The allocation simulation omits
shifts, travel, specialization, geography, and intra-day dynamics; crews
are interchangeable with fixed productivity $\kappa$. Conclusions are
therefore about the \emph{qualitative} structure of the
accuracy-to-decision mapping under swept parameters, not about any city's
achievable operational numbers. This is a simulation; nothing here is a
deployment.
\item \textbf{Reported demand only.} Targets and losses are functions of
reported requests; unreported need is invisible to both the forecasters
and the simulated planner (Section~10).
\item \textbf{Harmonization is a modeling choice.} The eight families are
deterministic keyword buckets; native taxonomies differ in granularity
and intake practice across cities, and family-level cross-city
comparisons inherit that heterogeneity. All mappings and compositions
are published for audit.
\item \textbf{Weather information set.} Target-day weather uses realized
observations as a proxy for a day-ahead forecast (A3); the
lagged-weather-only variant in Section~9 bounds the consequence. One
station per city coarsens spatial weather variation (A2).
\item \textbf{Four U.S. cities.} All systems are American Socrata
portals; generalization to differently structured municipal systems
(including the Gulf-region cities that motivate the broader research
programme) is untested and is future work, not an implied claim.
\item \textbf{No causal claims.} All results are predictive and
simulational; no intervention occurred in any city.
\end{itemize}
